\chapter{Effects, schedulers, and microtasks}
\label{ch:effects}

\begin{aside}
An effect is the workhorse of reactivity. The scheduler
decides when each effect runs. Microtasks are the small jobs
that need to happen ``soon'' but not synchronously.
\end{aside}

\section{Eager vs deferred}

An effect runs eagerly when its dependencies change. ``Eager''
in \maya{} means ``before the next frame''. A signal set
inside an event handler triggers effects after the handler
returns, before render.

\section{Order of execution}

When multiple effects are dirty, the runtime fires them in
topological order: an effect that depends on another's output
runs second.

\maya{} computes this order at registration time using the
dependency graph. Cycles, if any, break the order; the
runtime warns.

\section{Microtask queue}

A microtask is a deferred function that runs at the next
``safe point'' (between event handler and render). Used for:

\begin{itemize}[leftmargin=*]
  \item Continuations after async work.
  \item Cleanup that must happen before the next frame.
  \item Batched updates across multiple events.
\end{itemize}

Queue with \code{ctx.microtask(...)}; the loop drains
the queue before rendering.

\section{Idle callbacks}

For low-priority work (telemetry, prefetch), an idle
callback runs when the loop is otherwise quiet:

\begin{lstlisting}
ctx.on_idle([&] { precompute_thumbnail(); });
\end{lstlisting}

The callback fires after the loop has been idle for a
threshold (50 ms default).

\section{Effect strategies}

Different effect kinds:

\begin{itemize}[leftmargin=*]
  \item \code{effect}: runs on every dependency change.
  \item \code{effect\_throttled}: at most once per N ms.
  \item \code{effect\_debounced}: after N ms of quiet.
  \item \code{effect\_first}: runs once on first change.
\end{itemize}

\section{Cleanup}

An effect can register a cleanup to run before its next
invocation:

\begin{lstlisting}
ctx.effect([&] {
    auto timer = start_timer(...);
    return [timer] { cancel_timer(timer); };
});
\end{lstlisting}

The returned lambda is the cleanup. On the next run (or
unmount) the cleanup fires first.

\section{Async effects}

An effect can spawn an async job:

\begin{lstlisting}
ctx.effect([&] {
    auto id = ctx.get(user_id);
    spawn([id] {
        auto data = fetch(id);
        post_to_main_thread([&] {
            ctx.set(user_data, data);
        });
    });
});
\end{lstlisting}

The spawned thread does work; the main thread receives the
result via a post-back. Signal updates must happen on the
main thread.

\section{Suspension}

For coroutines, an effect can co\_await a future:

\begin{lstlisting}
ctx.async_effect([&]() -> Task {
    auto data = co_await fetch(...);
    ctx.set(loaded, data);
});
\end{lstlisting}

The runtime resumes the coroutine on the main thread when
the future completes.

\section{Recap}

\begin{recap}
\begin{itemize}[leftmargin=*]
  \item Effects fire eagerly between event and render.
  \item Topological order ensures dependents see updated
    inputs.
  \item Microtasks for ``soon, not now''.
  \item Idle callbacks for ``when nothing else needs to
    happen''.
  \item Throttle, debounce, cleanup, async variants.
\end{itemize}
\end{recap}

\section*{Workshop}

\begin{enumerate}[leftmargin=*]
  \item Build a debounced effect that saves text input
    after 500 ms of quiet.
  \item Use an idle callback for analytics.
  \item Implement an async effect with cleanup that
    cancels mid-flight.
\end{enumerate}
